\title{{{ ACL 2023} } Self-Instruct: Aligning Language Models with Self-Generated Instructions}

\begin{abstract}
Large ``instruction-tuned'' language models (i.e., finetuned to respond to instructions) have demonstrated a remarkable ability to generalize zero-shot to new tasks. Nevertheless, they depend heavily on human-written instruction data that is often limited in quantity, diversity, and creativity, therefore hindering the generality of the tuned model. We introduce \textsc{Self-Instruct}, a  framework for improving the instruction-following capabilities of pretrained language models by bootstrapping off their own generations. Our pipeline generates instructions, input, and output samples from a language model, then filters invalid or similar ones before using them to finetune the original model. Applying our method to the vanilla \textsc{GPT3}, we demonstrate a 33\% absolute improvement over the original model on \textsc{Super-NaturalInstructions}, on par with the performance of $\text{InstructGPT}_{\text{001}}$,\footnote{\label{footnote:gpt:instruct}Unless otherwise specified, our comparisons are with the {\tt text-davinci-001} engine. We focus on this engine since it is the closest to our experimental setup: supervised finetuning with human demonstrations. The newer engines are more powerful, though they use more data (e.g., code completion or latest user queries) or algorithms (e.g., PPO) that are difficult to compare with.} which was trained with private user data and human annotations. For further evaluation, we curate a set of expert-written instructions for novel tasks, and show through human evaluation that tuning GPT3 with \textsc{Self-Instruct} outperforms using existing public instruction datasets by a large margin, leaving only a 5\% absolute gap behind $\text{InstructGPT}_{\text{001}}$.
\textsc{Self-Instruct} provides an almost annotation-free method for aligning pretrained language models with instructions, and we release our large synthetic dataset to facilitate future studies on instruction tuning.\footnote{Code and data are available at \url{https://github.com/yizhongw/self-instruct}
}
\end{abstract}

\section{Introduction}
\label{sec:intro}

The recent  NLP literature has witnessed a tremendous amount of activity in building models that can follow natural language instructions~\citep[i.a.]{mishra2022cross,wei2022finetuned,sanh2022multitask,wang2022benchmarking,ouyang2022training,chung2022scaling}.
These developments are powered by two key components: large pretrained language models (LM) and human-written instruction data (e.g., 
\textsc{PromptSource}~\citep{bach2022promptsource} and \textsc{Super-NaturalInstructions}~\citep[\textsc{SuperNI} for short]{wang2022benchmarking}). However, collecting such instruction data is costly and often suffers limited diversity given that most human generations tend to be popular NLP tasks, falling short of covering a true variety of tasks and different ways to describe them. Continuing to improve the quality and coverage of instruction-tuned models necessitates the development of alternative approaches for supervising the instruction tuning process. 

In this work, we introduce \textsc{Self-Instruct}, a semi-automated process for instruction-tuning a pretrained LM using instructional signals from the model itself. The overall process is an iterative bootstrapping algorithm (see \autoref{fig:pipeline}), which starts off with a limited (e.g., 175 in our study) seed set of manually-written tasks that are used to guide the overall generation. In the first phase, the model is prompted to generate instructions for new tasks. This step leverages the existing collection of instructions to create more broad-coverage instructions that define (often new) tasks. Given the newly-generated set of instructions, the framework also creates input-output instances for them, which can be later used for supervising the instruction tuning. Finally, various heuristics are used to automatically filter low-quality or repeated instructions, before adding the remaining valid tasks to the task pool. This process can be repeated for many iterations until reaching a large number of tasks.

To evaluate \textsc{Self-Instruct} empirically, we run this framework on \textsc{GPT3}~\cite{brown2020gpt3}, which is a vanilla LM (\S\ref{sec:generatd-data}).
The iterative \textsc{Self-Instruct} process on this model leads to about 52k instructions, paired with about 82K instance inputs and target outputs. We observe that the resulting data provides a diverse range of creative tasks, as is demonstrated by examples in \autoref{fig:teaser-examples}. These generated tasks deviate from the distribution of typical NLP tasks, and also have fairly small overlap with the seed tasks (\S\ref{subsec:diversity}).
On this resulting data, we build \textsc{GPT3}$_{\textsc{Self-Inst}}$ by finetuning \textsc{GPT3} (i.e., the same model used for generating the instruction data). We evaluate \textsc{GPT3}$_{\textsc{Self-Inst}}$ in comparison to various other models on both typical NLP tasks included in \textsc{SuperNI}~\cite{wang2022benchmarking}, and a set of new instructions that are created for novel usage of instruction-following models (\S\ref{sec:experiment_results}). 
The results indicate that 
\textsc{GPT3}$_{\textsc{Self-Inst}}$ outperforms \textsc{GPT3} (the original model) by a large margin (+33.1\%) and nearly matches the performance of $\text{InstructGPT}_{\text{001}}$. Moreover, our human evaluation on the newly-created instruction set shows that \textsc{GPT3}$_{\textsc{Self-Inst}}$ demonstrates a broad range of instruction following ability, outperforming models trained on other publicly available instruction datasets and leaving only a 5\% gap behind $\text{InstructGPT}_{\text{001}}$. 

In summary, our contributions are: (1) we introduce \textsc{Self-Instruct}, a method for inducing instruction following capabilities with minimal human-labeled data; (2) we demonstrate its effectiveness via extensive instruction-tuning experiments; and (3) we release a large synthetic dataset of 52K instructions and a set of manually-written novel tasks for building and evaluating future instruction-following models.

\section{Method}
\label{sec:method}

Annotating large-scale instruction data can be challenging for humans because it requires 1) creativity to come up with novel tasks and 2) expertise for writing the solutions to each task. Here, we detail our process for \textsc{Self-Instruct}, which refers to the pipeline of generating tasks with a \textit{vanilla pretrained language model} itself, filtering the generated data, and then conducting instruction tuning with this generated data in order to align the LM to follow instructions better. This pipeline is depicted in \autoref{fig:pipeline}. 

\subsection{Defining Instruction Data}
\label{subsec:instruction-data-definition}

The instruction data we want to generate contains a set of instructions $\{I_t\}$, each of which defines a task $t$ in natural language. Task $t$ has $n_t \ge 1$ input-output instances $\{(X_{t,i}, Y_{t,i})\}_{i=1}^{n_t}$.
A model $M$ is expected to produce the output, given the task instruction and the corresponding input: $M(I_t, X_{t,i}) = Y_{t,i}$, for $i \in \{1,\ldots,n_t\}$.
Note that the instruction and instance input does not have a strict boundary in many cases. For example, ``write an essay about school safety'' can be a valid instruction that we expect models to respond to directly, while it can also be formulated as ``write an essay about the following topic'' as the instruction, and ``school safety'' as an instance input. To encourage the diversity of the data format, we allow such instructions that do not require additional input (i.e., $X$ is empty).

\subsection{Automatic Instruction Data Generation}
\label{subsec:data-generation}

Our pipeline for data generation consists of four steps: 1) generating task instructions, 2) determining if the instruction represents a classification task, 3) instance generation with either an input-first or output-first approach, and 4) filtering low-quality data.

\paragraph{Instruction Generation.} At the first step, \textsc{Self-Instruct} generates new instructions from a small set of seed human-written instructions in a bootstrapping fashion. We initiate the task pool with 175 tasks (1 instruction and 1 instance for each task).\footnote{These tasks were newly written by the authors and their labmates at UW, without reference to existing datasets or the test set used in this work. We provide more details about these tasks and analyze their similarity to the test tasks in Appendix \S\ref{subsec:seed_tasks}.} For every step, we sample 8 task instructions from this pool as in-context examples. Of the 8 instructions, 6 are from the human-written tasks, and 2 are from the model-generated tasks in previous steps to promote diversity. The prompting template is shown in \autoref{tab:instruction_generation_template}.

\paragraph{Classification Task Identification.} Because we need two different approaches for classification and non-classification tasks, we next identify whether the generated instruction represents a classification task or not.\footnote{More concretely, we regard tasks that have a small limited output label space as classification tasks.} We prompt the LM in a few-shot way to determine this, using 12 classification instructions and 19 non-classification instructions from the seed tasks. The prompting template is shown in \autoref{tab:classification_task_identification_template}. 

\paragraph{Instance Generation.} Given the instructions and their task type, we generate instances for each instruction independently. This is challenging because it requires the model to understand what the target task is, based on the instruction, figure out what additional input fields are needed and generate them, and finally complete the task by producing the output. We found that pretrained LMs can achieve this to a large extent when prompted with instruction-input-output in-context examples from other tasks. A natural way to do this is the \textbf{Input-first Approach}, where we can ask an LM to come up with the input fields first based on the instruction, and then produce the corresponding output. This generation order is similar to how models are used to respond to instruction and input, but here with in-context examples from other tasks. The prompting template is shown in \autoref{tab:input-first-generation-template}. 

However, we found that this approach can generate inputs biased toward one label, especially for classification tasks (e.g., for grammar error detection, it usually generates grammatical input). Therefore, we additionally propose an
\textbf{Output-first Approach} for classification tasks, where we first generate the possible class labels, and then condition the input generation on each class label. The prompting template is shown in \autoref{tab:output-first-generation-template}.\footnote{In this work, we use a fixed set of seed tasks for prompting the instance generation, and thus only generate a small number of instances per task in one round. Future work can use randomly sampled tasks to prompt the model to generate a larger number of instances in multiple rounds.} We apply the output-first approach to the classification tasks identified in the former step, and the input-first approach to the remaining non-classification tasks.

\paragraph{Filtering and Postprocessing.}
\label{subsec:filtering}
To encourage diversity, a new instruction is added to the task pool only when its ROUGE-L similarity with any existing instruction is less than 0.7. We also exclude instructions that contain some specific keywords (e.g., image, picture, graph) that usually can not be processed by LMs. When generating new instances for each instruction, we filter out instances that are exactly the same or those with the same input but different outputs. Invalid generations are identified and filtered out based on heuristics (e.g., instruction is too long or too short, instance output is a repetition of the input). 

\subsection{Finetuning the LM to Follow Instructions}
\label{subsec:instruction-tuning}
After creating large-scale instruction data, we use it to finetune the original LM (i.e., \textsc{Self-Instruct}). To do this, we concatenate the instruction and instance input as a prompt and train the model to generate the instance output in a standard supervised way. To make the model robust to different formats, we use multiple templates to encode the instruction and instance input together. For example, the instruction can be prefixed with ``Task:'' or not, the input can be prefixed with ``Input:'' or not, ``Output:'' can be appended at the end of the prompt or not, and different numbers of break lines can be put in the middle, etc.

\section{\textsc{Self-Instruct} Data from \textsc{GPT3}}
\label{sec:generatd-data}

 In this section, we apply our method for inducing instruction data to \textsc{GPT3} as a case study. We use the largest GPT3 LM (``davinci'' engine) accessed through the OpenAI API.\footnote{\url{https://openai.com/api/}} The parameters for making queries are described in Appendix \ref{subsec:query_gpt3_api}. Here we present an overview of the generated data.

\subsection{Statistics}
\label{subsec:statistics}

\autoref{tab:data_statistics} describes the basic statistics of the generated data. We generate a total of over 52K instructions and more than 82K instances corresponding to these instructions after filtering.

\subsection{Diversity}
\label{subsec:diversity}
To study what types of instructions are generated and how diverse they are, we identify the verb-noun structure in the generated instructions. We use the Berkeley Neural Parser\footnote{\url{https://parser.kitaev.io/}} \citep{kitaev-klein-2018-constituency,kitaev-etal-2019-multilingual} to parse the instructions and then extract the verb that is closest to the root as well as its first direct noun object. 26,559 out of the 52,445 instructions contain such structure; other instructions usually contain more complex clauses (e.g., ``Classify whether this tweet contains political content or not.'') or are framed as questions (e.g., ``Which of these statements are true?''). We plot the top 20 most common root verbs and their top 4 direct noun objects in \autoref{fig:verb-noun-distribution}, which account for 14\% of the entire set. Overall, we see quite diverse intents and textual formats in these instructions.

We further study how the generated instructions differ from the seed instructions used to prompt the generation. For each generated instruction, we compute its highest ROUGE-L overlap with the 175 seed instructions. We plot the distribution of these ROUGE-L scores in \autoref{fig:overlap-distribution}. The results indicate a decent number of new instructions were generated, which do not have much overlap with the seeds. We also demonstrate diversity in the length of the instructions, instance inputs, and instance outputs in \autoref{fig:length_distribution}. 

\subsection{Quality}
So far, we have shown the quantity and diversity of the generated data, but its quality remains uncertain. To investigate this, we randomly sample 200 instructions and randomly select 1 instance per instruction. We asked an expert annotator (author of this work) to label whether each instance is correct or not, in terms of the instruction, the instance input, and the instance output. Evaluation results in \autoref{tab:data_quality_eval} show that most of the generated instructions are meaningful, while the generated instances may contain more noise (to a reasonable extent). However, we found that even though the generations may contain errors, most of them are still in the correct format or partially correct, which can provide useful guidance for training models to follow instructions. We listed a number of good examples and bad examples in \autoref{tab:generated_tasks} and \ref{tab:bad-generated-tasks}, respectively.

\label{subsec:quality}

\section{Experimental Results}
\label{sec:experiment_results}

We conduct experiments to measure and compare the performance of models under various instruction tuning setups. We first describe our models and other baselines, followed by our experiments. 

\subsection{\textsc{GPT3}$_{\textsc{Self-Inst}}$: finetuning \textsc{GPT3} on its own instruction data}
\label{subsec:self-instruct-gpt3}
Given the instruction-generated instruction data, we conduct instruction tuning with the \textsc{GPT3} model itself (``davinci'' engine). As described in \S\ref{subsec:instruction-tuning}, we use various templates to concatenate the instruction and input, and train the model to generate the output. This finetuning is done through the OpenAI finetuning API.\footnote{See \href{https://beta.openai.com/docs/guides/fine-tuning}{OpenAI's documentation on finetuning}. } We use the default hyper-parameters, except that we set the prompt loss weight to 0, and we train the model for 2 epochs. We refer the reader to Appendix~\ref{subsec:finetuning_details} for additional finetuning details. The resulting model is denoted by \textsc{GPT3}$_{\textsc{Self-Inst}}$.

\subsection{Baselines}
\paragraph{Off-the-shelf LMs.} We evaluate T5-LM \cite{lester2021power,raffel2020exploring} and \textsc{GPT3}~\cite{brown2020gpt3} as the vanilla LM baselines (only pretraining, no additional finetuning). These baselines will indicate the extent to which off-the-shelf LMs are capable of following instructions naturally immediately after pretraining.

\paragraph{Publicly available instruction-tuned models.}
\textsc{T$0$} and \textsc{T$k$-Instruct} are two instruction-tuned models proposed in \citet{sanh2022multitask} and \citet{wang2022benchmarking}, respectively, and are demonstrated to be able to follow instructions for many NLP tasks. Both of these models are finetuned from the T5~\cite{raffel2020exploring} checkpoints and are publicly available.\footnote{
\textsc{T$0$} is available at \href{https://huggingface.co/bigscience/T0}{here} and \textsc{T$k$-Instruct} is  \href{https://huggingface.co/allenai/tk-instruct-11b-def}{here}.
} For both of these models, we use their largest version with 11B parameters.

\paragraph{Instruction-tuned GPT3 models.} We evaluate $\text{InstructGPT}_{\text{}}$~\cite{ouyang2022training}, 
which is developed by OpenAI based on GPT3 to follow human instructions better and has been found by the community to have impressive zero-shot abilities. There are various generations of these models, where newer ones use more expansive data or algorithmic novelties.\footnote{
See \href{https://beta.openai.com/docs/model-index-for-researchers}{OpenAI's documentation on their models.} } For our \textsc{SuperNI} experiments in \S\ref{subsec:superni-experiments}, we only compare with their \texttt{text-davinci-001} engine, because their newer engines are trained with the latest user data and are likely to have already seen the \textsc{SuperNI} test set. For our human evaluation on newly written instructions, we include their 001, 002 and 003 engines for completeness.

Additionally, to compare \textsc{Self-Instruct} training with other publicly available instruction tuning data, we further finetune GPT3 model with data from \textsc{PromptSource} and \textsc{SuperNI}, which are used to train the \textsc{T$0$} and \textsc{T$k$-Instruct} models. We call them \textsc{T$0$} training and \textsc{SuperNI} training for short, respectively. To save the training budget, we sampled 50K instances (but covering all their instructions) for each dataset, which has a comparable size to the instruction data we generated. Based on the findings from \citet{wang2022benchmarking} and our early experiments, reducing the number of instances per training task does not degrade the model's generalization performance to unseen tasks.

\subsection{Experiment 1: Zero-Shot Generalization on \textsc{SuperNI} benchmark}
\label{subsec:superni-experiments}
We first evaluate the models' ability to follow instructions on typical NLP tasks in a zero-shot fashion. We use the evaluation set of \textsc{SuperNI} ~\cite{wang2022benchmarking}, which consists of 119 tasks with 100 instances in each task. In this work, we mainly focus on the zero-shot setup, i.e., the model is prompted with the definition of the tasks only, without in-context demonstration examples. For all our requests to the \textsc{GPT3} variants, we use the deterministic generation mode (temperature as 0 and no nucleus sampling) without specific stop sequences.

\paragraph{Results.}
We make the following observations from the results in \autoref{tab:superni_results}. 
\textsc{Self-Instruct} boosts the instruction-following ability of \textsc{GPT3} by a large margin. The vanilla \textsc{GPT3} model basically cannot follow human instructions at all. Upon manual analysis, we find that it usually generates irrelevant and repetitive text, and does not know when to stop generation. Compared with other models that are not specifically trained for \textsc{SuperNI}, \textsc{GPT3}$_{\textsc{Self-Inst}}$ achieves better performance than \textsc{T$0$} or the \textsc{GPT3} finetuned on the \textsc{T$0$} training set, which takes tremendous human labeling efforts. Notably, \textsc{GPT3}$_{\textsc{Self-Inst}}$ also nearly matches the performance of $\text{InstructGPT}_{\text{001}}$, which is trained with private user data and human-annotated labels.

Models trained on the \textsc{SuperNI} training set still achieve better performance on its evaluation set, which we attribute to the similar instruction style and formatting. However, we show that 
\textsc{Self-Instruct} still brings in additional gains when combined with the \textsc{SuperNI} training set, proving its value as complementary data.

\begin{comment}
\newcommand{-70ex}{-70ex}

\end{comment}

\newcommand{-54.5ex}{-54.5ex}
\newcommand{2ex}{2ex}

\subsection{Experiment 2: Generalization to User-oriented Instructions on Novel Tasks}
\label{sec:user_instructions}

Despite the comprehensiveness of \textsc{SuperNI} in collecting existing NLP tasks, most of these NLP tasks were proposed for research purposes and skewed toward classification. To better access the practical value of instruction-following models, a subset of the authors curate a new set of instructions motivated by user-oriented applications. We first brainstorm various domains where large LMs may be useful (e.g., email writing, social media, productivity tools, entertainment, programming), then craft instructions related to each domain along with an input-output instance (again, input is optional). We aim to diversify the styles and formats of these tasks (e.g., instructions may be long or short; input/output may take the form of bullet points, tables, codes, equations, etc.). In total, we create 252 instructions with 1 instance per instruction. We believe it can serve as a testbed for evaluating how instruction-based models handle diverse and unfamiliar instructions. \autoref{tab:case_study} presents a small portion of them. The entire set is available in our GitHub repository. We analyze the overlap between this set set and the seed instructions in \S\ref{subsec:seed_tasks}.

\paragraph{Human evaluation setup.} Evaluating models' performance on this evaluation set of diverse tasks is extremely challenging because different tasks require different expertise. Indeed, many of these tasks cannot be measured by automatic metrics or even be judged by normal crowdworkers (e.g., writing a program, or converting first-order logic into natural language). To get a more faithful evaluation, we asked the authors of the instructions to judge model predictions. Details on how we set up this human evaluation are described in Appendix \ref{sec:human_eval_details}. The evaluators were asked to rate the output based on whether it accurately and effectively completes the task. We implemented a four-level rating system for categorizing the quality of the models' outputs:

\begin{itemize}[noitemsep, leftmargin=*]
    \item \textsc{Rating-A:} The response is valid and satisfying.
    \item \textsc{Rating-B:} The response is acceptable but has minor errors or imperfections.
    \item \textsc{Rating-C:} The response is relevant and responds to the instruction, but it has significant errors in the content. For example, GPT3 might generate a valid output first, but continue to generate other irrelevant things.
    \item \textsc{Rating-D:} The response is irrelevant or completely invalid.
\end{itemize}

\paragraph{Results.}
\autoref{fig:human_instructions_results} shows the performance of \textsc{GPT3} model and its instruction-tuned counterparts on this newly written instruction set (w. inter-rater agreement  $\kappa=0.57$ on the 4-class categorical scale, see Appendix \ref{sec:human_eval_details} for details). As anticipated, the vanilla \textsc{GPT3} LM is largely unable to respond to instructions, and all instruction-tuned models demonstrate comparatively higher performance. Nonetheless, \textsc{GPT3}$_{\textsc{Self-Inst}}$ (i.e.,  \textsc{GPT3} model finetuned with \textsc{Self-Instruct}) outperforms those counterparts trained on \textsc{T$0$} or \textsc{SuperNI} data by a large margin, demonstrating the value of the generated data despite the noise. Compared with $\text{InstructGPT}_{\text{001}}$, \textsc{GPT3}$_{\textsc{Self-Inst}}$ is quite close in performance---if we count acceptable response with minor imperfections (\textsc{Rating-B}) as valid, \textsc{GPT3}$_{\textsc{Self-Inst}}$ is only 5\% behind $\text{InstructGPT}_{\text{001}}$. Lastly, our evaluation confirms the impressive instruction-following ability of $\text{InstructGPT}_{\text{002}}$ and $\text{InstructGPT}_{\text{003}}$. Although there are many factors behind this success, we conjecture that future work can largely benefit from improving the quality of our generated data by using human annotators or training a reward model to select better generations, similar to the algorithm used by \citet{ouyang2022training}.

\subsection{Effect of Data Size and Quality}
\label{subsec:size-and-quality-analysis}

\paragraph{Data size.}
\textsc{Self-Instruct} provides a way to grow instruction data at a low cost with almost no human labeling; could more of this generated data lead to better instruction-following ability? We conduct an analysis of the size of generated data by subsampling different numbers of instructions from the generated dataset, finetuning \textsc{GPT3} on the sampled subsets, and evaluating how the resulting models perform on the 252 user-oriented instruction set. We conduct the same human evaluation as in \S\ref{sec:user_instructions}.
\autoref{fig:data_scaling_analysis} presents the performance of \textsc{GPT3}$_{\textsc{Self-Inst}}$ models finetuned with different sizes of generated data. Overall, we see consistent improvement as we grow the data size. However, this improvement almost plateaus after 16K. This is in-line with the data scaling experiments in \citet[Fig.~5]{wang2022benchmarking}.
Interestingly, when evaluating on \textsc{SuperNI} we found the model's performance gain plateaus earlier at around hundreds of instructions. This may be due to the fact that the new generated data is distinct from typical NLP tasks in \textsc{SuperNI}, indicating that future research may benefit from using a combination of different instruction data for better performance on various types of tasks.

\paragraph{Data quality.} Another direction to improve the model's performance is to take our generated data and get better supervision (with less noise). We explore this idea by using $\text{InstructGPT}_{\text{003}}$ (the best available general-purpose model) to regenerate the output field of all our instances given the instruction and input. We then use this improved version of our data to finetune \textsc{GPT3}. This can be regarded as a distillation of $\text{InstructGPT}_{\text{003}}$ with our data. As is shown in \autoref{fig:data_scaling_analysis}, the resulting model outperforms the counterpart trained with the original data by 10\%, which suggests big room for future work on using our generation pipeline to get initial data and then improving the data quality with human experts or distillation from better models.

\section{Related Work}
\label{sec:related}

\paragraph{Instruction-following LMs.} A series of works have found evidence that vanilla LMs can be effective at following general language instructions if tuned with annotated ``instructional'' data---datasets containing language instructional commands and their desired outcomes based on human annotation~\cite[i.a.]{weller2020learning, mishra2022cross,wei2022finetuned,sanh2022multitask}. Additionally, they show a direct correlation between the size and diversity of the ``instructional'' data and the generalizability of resulting models to unseen tasks \citep{wang2022benchmarking,chung2022scaling}. However, since these developments largely focus on existing NLP tasks and depend on human-annotated instructions, this poses a bottleneck for progress toward more generalizable models 
\cite[e.g., see Fig.~5a in][]{wang2022benchmarking}. Our work aims to move beyond classical NLP tasks and tackle the challenges of creating diverse instruction data by employing pretrained LMs.
$\text{InstructGPT}_{\text{}}$~\cite{ouyang2022training} shares a similar goal as ours in building more general-purpose LMs, and has demonstrated remarkable performance in following diverse user instructions. However, as a commercial system, their construction process still remains quite opaque. In particular, the role of \emph{data} has remained understudied due to limited transparency and the private user data they used in their study.

Addressing such challenges necessitates the creation of a large-scale, public dataset covering a broad range of tasks.

\paragraph{Language models for data generation and augmentation.} A variety of works have proposed using LMs for data generation~\cite{schick2021generating,wang2021towards,liu2022wanli,meng2022tuning} or augmentation~\cite{feng2021survey,yang2020generative,mekala2022intermediate}. Our work differs from this line in that it is \emph{not} specific to a particular task (say, QA or NLI). In contrast, a distinct motivation for \textsc{Self-Instruct} is to bootstrap new task definitions that may not have been defined before by NLP practitioners (though potentially still important for real users).  In parallel with our work, \citet{honovich2022unnatural} also propose to generate large-scale instruction data (so-called Unnatural Instructions) with GPT3 models. The major differences are that 1) they use tasks in \textsc{SuperNI}~\citep{wang2022benchmarking} as their seed tasks, resulting in a different distribution of generated tasks; 2) they employ $\text{InstructGPT}_{\text{002}}$ for generating the data, in which sense they are distilling knowledge from an already instruction-tuned model, while we solely rely on the vanilla LM; 3) the detailed generation pipeline and templates are different. Nevertheless, we believe that both efforts in expanding instruction data are complementary, and the community will benefit from these diverse datasets.

\paragraph{Instruction generation.} A series of recent works~\cite{zhou2022large, ye2022guess, singh2022explaining, honovich2022instruction} generate instructions of a task given a few examples. While \textsc{Self-Instruct} also involves instruction generation, a major difference in our case is it is task-agnostic; we generate new tasks (instructions along with instances) from scratch.

\paragraph{Model self-training.} A typical self-training framework ~\cite{he2019revisiting, xie2020self, du2021self, amini2022self, huang2022large} uses trained models to assign labels to unlabeled data and then leverages the newly labeled data to improve the model. In a similar line, \citet {zhou2022prompt} use multiple prompts to specify a single task and propose to regularize via prompt consistency, encouraging consistent predictions over the prompts. This allows either finetuning the model  with extra unlabeled training data, or direct application at inference time. While \textsc{Self-Instruct} has similarities with the self-training literature, most self-training methods assume a specific \textit{target task} as well as \textit{unlabeled examples} under it; in contrast, \textsc{Self-Instruct} produces a variety of tasks from scratch.

\paragraph{Knowledge distillation.} Knowledge distillation~\cite{hinton2015distilling,Sanh2019DistilBERTAD, west2021symbolic,luciecharlotte2022teachingsmallmodels} often involves the transfer of knowledge from larger models to smaller ones. \textsc{Self-Instruct} can also be viewed as a form of ``knowledge distillation", however, it differs from this line in the following ways: (1) the source and target of distillation are the same, i.e., a model's knowledge is distilled to itself; (2) the content of distillation is in the form of an instruction task (i.e., instructions that define a task, and a set of examples that instantiate it).

\paragraph{Bootstrapping with limited resources.} A series of recent works use language models to bootstrap some inferences using specialized methods. NPPrompt~\cite{zhao2022pre} provides a method to generate predictions for semantic labels without any finetuning. It uses a model's own embeddings to automatically find words relevant to the label of the data sample and hence reduces the dependency on manual mapping from model prediction to label (verbalizers). STAR~\cite{zelikman2022star} iteratively leverages a small number of rationale examples and a large dataset without rationales, to bootstrap a model's ability to perform reasoning.  Self-Correction~\cite{welleck2022generating} decouples an imperfect base generator (model) from a separate corrector that learns to iteratively correct imperfect generations and demonstrates improvement over the base generator. Our work instead focuses on bootstrapping new tasks in the instruction paradigm.

\paragraph{Multi-modal instruction-following.} Instruction-following models have also been of interest in the multi-modal learning literature \cite{fried2018speaker, shridhar2020alfred, min2022film, weir2022one}.
\textsc{Self-Instruct}, as a general approach to expanding data, can potentially also be helpful in those settings, which we leave to future work.

\section{Conclusion}
\label{sec:conclusion}
We introduce \textsc{Self-Instruct}, a method to improve the instruction-following ability of LMs via their own generation of instruction data. On experimenting with vanilla \textsc{GPT3}, we automatically construct a large-scale dataset of 52K instructions for diverse tasks, and finetuning GPT3 on this data leads to a 33\% absolute improvement on \textsc{SuperNI} over the original \textsc{GPT3}. Furthermore, we curate a set of expert-written instructions for novel tasks. Human evaluation on this set shows that tuning GPT3 with \textsc{Self-Instruct} outperforms using existing public instruction datasets by a large margin and performs closely to $\text{InstructGPT}_{\text{001}}$.
We hope \textsc{Self-Instruct} can serve as the first step to align pretrained LMs to follow human instructions, and future work can build on top of this data to improve instruction-following models.

\section{Broader Impact}

Beyond the immediate focus of this paper, we believe that \textsc{Self-Instruct} may help bring more transparency to what happens ``behind the scenes'' of widely-used instruction-tuned models like $\text{InstructGPT}_{\text{}}$ or ChatGPT. Unfortunately, such industrial models remain behind API walls as their datasets are not released, and hence there is little understanding of their construction and why they demonstrate impressive capabilities. The burden now falls on academia to better understand the source of success in these models and strive for better---and more open---models. We believe our findings in this paper demonstrate the importance of diverse instruction data, and our large synthetic dataset can be the first step toward higher-quality data for building better instruction-following models. 

At this writing, the central idea of this paper has been adopted in several follow-up works for such endeavors \citep[i.a.]{alpaca,xu2023baize,sun2023principle}. 

\section{Limitations}
Here, we discuss some limitations of this work to inspire future research in this direction.  

\paragraph{Tail phenomena.} 
\textsc{Self-Instruct} depends on LMs, and it will inherit all the limitations that carry over with LMs. As recent studies have shown~\cite{razeghi2022impact,kandpal2022large},  \emph{tail phenomena} pose a serious challenge to the success of LMs. In other words, LMs' largest gains correspond to the frequent uses of languages (head of the language use distribution), and there might be minimal gains in the low-frequency contexts. Similarly, in the context of this work, it would not be surprising if the majority of the gains by \textsc{Self-Instruct} are skewed toward tasks or instructions that present more frequently in the pretraining corpus. As a consequence, the approach might show brittleness with respect to uncommon and creative instructions.

\paragraph{Dependence on large models.} Because of \textsc{Self-Instruct}'s dependence on the inductive biases extracted from LMs, it might work best for larger models. If true, this may create barriers to access for those who may not have large computing resources. We hope future studies will carefully study the gains as a function of model size or various other parameters. It is worthwhile to note that instruction-tuning with human annotation also suffers from a similar limitation: gains of instruction-tuning are higher for larger models~\cite{wei2022finetuned}. 

\paragraph{Reinforcing LM biases.} A point of concern for the authors is the unintended consequences of this iterative algorithm, such as the amplification of problematic social biases (stereotypes or slurs about gender, race, etc.). Relatedly, one observed challenge in this process is the algorithm's difficulty in producing balanced labels, which reflected models' prior biases. We hope future work will lead to better understanding of the pros and cons of the approach. 

\end{document}
